\documentclass[10pt]{article}
\usepackage[letterpaper,landscape,margin=9mm]{geometry}
\usepackage{graphicx}
\usepackage{caption}
\pagestyle{empty}
\graphicspath{{output/pilot_cross_dataset/}}
\captionsetup{font=small,labelfont=bf}
\begin{document}
\begin{figure}[p]
\centering
\includegraphics[width=0.94\textwidth]{comparison_existing_3d.pdf}
\caption*{\textbf{Pilot: terminal parentage across molecular representations,
existing 3D coordinate convention.} All curves use the same 188 natural
terminal parent--child edges. C. elegans embryo 1 uses cutoff 255; the two
AF16 C. briggsae embryos use cutoffs 148 and 156. Protein uses the frozen
top-20 z-scored features; RNA uses the frozen shared 20 TFs with cosine
distance on stored values. Panel A normalizes each objective by its own
endpoint span; crosses mark natural lineages and circles their nearest
sampled assignments. Panel B shows biological-parent retention along
canonical front position. The first-cousin null supplies exact moments;
$d_{NP}$ uses its mean and the same assignment selected by $d_{LP}$.
Endpoints minimize the other objective within numerical primary-optimum
tolerance. Lines connect sampled solutions, and retention maxima are observed
maxima. \textbf{C. briggsae axial calibration is unresolved; this is not a
publication-ready species comparison.}}
\end{figure}
\clearpage
\begin{figure}[p]
\centering
\includegraphics[width=0.94\textwidth]{comparison_xy.pdf}
\caption*{\textbf{Pilot sensitivity: the same terminal comparison using XY
positions only.} Molecular matrices, natural edges, parent-slot
multiplicities, and feature panels are held fixed. Only the travel matrix
changes, after which null moments, the assignment sweep, endpoints, and
canonical coordinates are recomputed. The smaller retention differences and
changed proximity ordering show sensitivity to spatial geometry. XY
projection does not establish a calibrated 3D species effect. The analysis
uses 301 weights; a nested 1,201-weight sweep preserves the reported proximity
and maximum-retention values to numerical precision. The matched set retains
only 1 of 32 programmed-death cells in the original protein analysis;
coverage and full-set comparisons are included in the separate four-page
\texttt{pilot\_review.pdf}. Both tracking replicates use the same C. briggsae
RNA matrix, so their spread does not estimate RNA measurement uncertainty.}
\end{figure}
\end{document}
